\documentclass[11pt]{article}
\usepackage[T1]{fontenc}
\usepackage{lmodern,amsmath,amssymb,amsthm,booktabs,geometry,hyperref}
\geometry{margin=1in}
\hypersetup{colorlinks=true,urlcolor=blue}
\newtheorem{theorem}{Theorem}
\newtheorem{lemma}{Lemma}
\newcommand{\DB}{D_B}
\newcommand{\Tr}{T}
\title{Bulgarian solitaire: progress on $\DB(n)$}
\author{P3 C6 research progress}
\date{26 September 2026}
\begin{document}
\maketitle

\begin{abstract}
This is a partial submission to P3 C6, which asks for $\DB(n)$ for every $n$.
The general problem remains open in this work.  A written argument gives the
small-residue formula at residue $3$ for all sufficiently large ranks; exhaustive
computation covers the remaining ranks.  Additional finite values are obtained
by exhaustive enumeration and a case-certificate checker.  We distinguish
these results from conjectural inverse-height bounds needed for a general proof.
\end{abstract}

\section{Problem and main result}
Let $B$ remove one card from each pile of a partition, discard empty piles,
create one pile equal to the old pile count, and sort.  Put
\[
 d_B(\lambda)=\min\{t\geq0:B^t(\lambda)\text{ is periodic}\},\qquad
 \DB(n)=\max_{\lambda\vdash n}d_B(\lambda),\qquad
 \Tr_j=\frac{j(j+1)}2.
\]
For $n=\Tr_{k-1}+3$, the following is the result supported by the research
report and accompanying exact computations.

\begin{theorem}[Residue $3$]
\[
\begin{array}{c|rrrrrr|c}
k&3&4&5&6&7&8&k\geq9\\ \hline
\DB(\Tr_{k-1}+3)&6&7&9&13&18&24&(k-1)(k-5)
\end{array}
\]
For $k\geq14$ there is a written upper-bound proof.  For $3\leq k\leq14$
the stated values have exhaustive backward-layer certificates, with $k=14$
covered by both methods.
\end{theorem}

The result uses the cyclic-partition classification and pile-count sandwich
lemmas of Griggs--Ho~\cite{GH}, also proved in the earlier P3 C1--C3 work.
Those prerequisites were read and independently stress-tested here; they are
not newly formalized in Lean.  The lower-bound construction is from
Griggs--Ho, with its trajectory calculation written out in the research
report.  No claim of first discovery or of a full C6 solution is made.

\section{Lower bound and proof structure}
For $k\geq9$ set $F=(k-1)(k-5)$ and
\[
 \lambda_k=(k-2,k-2,k-3,\ldots,5,4,4,3,2,1)\vdash\Tr_{k-1}+3.
\]
Relative to the staircase $(k-1,k-2,\ldots,1,0)$, this partition has a hole
on diagonal $k-1$ at column $0$ and four extra cells on diagonal $k$ at
columns $k-4,k-3,k-2,k-1$.  For $0\leq s<F$, write
$s=a(k-1)+b$, where $0\leq a\leq k-6$ and $0\leq b\leq k-2$.
The hole is at column $b$; the four extra cells are at
$(b-a-1-j)\bmod k$, $0\leq j\leq3$.  None meets the hole or its right
neighbor, so the unsorted move remains a partition.  The hole persists until
step $F$, when the first sort yields the cyclic boundary
$(k-1,k-2,k-2,k-3,k-4,k-6,\ldots,1)$.  Thus
$d_B(\lambda_k)=F$.  This calculation reproduces the lower-bound family of
Griggs--Ho Theorem 4.5, case (1).

For the upper bound, write $c_i$ for pile count at time $i$.  The cited
final-pattern dichotomy says a sufficiently late first cyclic time has an
earlier type-I sandwich $(k-1,k,\ldots,k,k+1)$ or type-II sandwich
$(k-2,k-1,\ldots,k-1,k)$.  For a sandwich
$(x-1,x,\ldots,x,x+1)$ of width $L$, its start time $p$ satisfies
$p\leq x(L-1)$.  At time $p+1$, the newborn pile has size $x-1$; among old
piles there is exactly one of each size $1,\ldots,L-2$, none of size
$L-1$, and all remaining old piles have size at least $L$.  This entry-state
description and card budget limit possible widths and states.

For $k\geq14$, the budget rules out type-I width $k-5$ or greater; retreat
then gives $d_B\leq F-6$.  Type-II widths at most $k-5$ give
$d_B\leq F-1$, and width $k-1$ is impossible.  Widths
$k-4,k-3,k-2,k$ require entry-state analysis.  At width $k-4$, the excess
over minimum pile sizes is $6$, spread over four free piles.  Its nine
partitions give four cyclic entry states, two states controlled by a
dominant-pile inverse-height bound, one state exposing an earlier short
sandwich, and two states excluded by a persistent diagonal-$(k+1)$ cell.
The cyclic state entered by $\lambda_k$ is the tight case.  At widths
$k-3$ and $k-2$, explicit inverse branches of the cyclic entry states end
in a dominant-pile state, a dead end, or an earlier short sandwich.  Full
width $k$ has a single possible entry state.  All resulting bounds are at
most $F$.

Two reusable bounds make this finite case analysis possible.  If a partition
$\mu\vdash n$ has a unique largest part $M$ exceeding every other part by
at least $3$, its inverse-tree height satisfies $h(\mu)\leq n-M+1$.
If its two largest parts are $M\geq M'\geq3$ and all others are at most
$M'-3$, then
\[
 h(\mu)\leq\left\lfloor\frac{n-M-M'}2\right\rfloor+2.
\]
For the first bound, every inverse step selecting a smaller part raises the
dominant part by one, so there are at most $n-M$ such steps; selecting the
dominant part ends the chain.  For the second, each tail step raises both
dominant parts and at most two head steps follow.  The research report lists
every residual entry state and inverse branch; all symbolic identities were
tested for $9\leq k\leq60$.  These finite checks support the written proof
but do not replace its all-$k$ reasoning.

\section{Exact finite evidence beyond residue $3$}
Backward layering from all cyclic partitions visits each partition exactly
once because every partition has one forward successor.  Summing layer sizes
and checking against the partition number $p(n)$ certifies exhaustive
coverage.  The computation agrees with the Griggs--Ho conjectured value
$L(n)$ for every $n\leq70$, and also for $n=81,82,94,95$.
At residue $3$, the new finite checks include
\[
 \DB(58)=60,\quad \DB(69)=77,\quad
 \DB(81)=96,\quad \DB(94)=117.
\]
The $n=94$ run covers $92{,}669{,}720$ partitions.  A separate case-certificate
checker applies the proved sandwich and inverse-height rules at fixed
$(r,k)$ to obtain the following additional exact values:
\[
\begin{array}{c|c|c}
r& k&\DB(\Tr_{k-1}+r)\\ \hline
4&12,\ldots,21&(k-1)(k-6)\\
5&15,\ldots,21&(k-1)(k-7)\\
6&17,\ldots,20&(k-1)(k-8)\\
7&19,\ldots,21&(k-1)(k-9)
\end{array}
\]
These are 24 fixed-$n$ computer-assisted results, not a theorem for every
rank.  Three overlap exhaustive enumeration and agree.  The checker has not
received an independent audit; its rule tags and replay code are retained
with the research output.

\section{Open step toward C6}
For general residue $r$ at type-II width $k-1-r$, the card excess is
$\Tr_r$ among $r+1$ free old piles, while the retreat bound equals
$F_r=(k-1)(k-r-2)$.  This identifies a tight case.  Other entry states
at larger widths pose an inverse-height question.  The proposed bound
\[
 \mu_1\geq m+3\quad\Longrightarrow\quad h(\mu)\leq n-\mu_1+1
 \qquad(\mu\vdash n,\;m=\text{number of parts})
\]
has been tested on every partition of $n\leq34$ but remains \emph{unproved}.
It follows from two still-unproved four-ones and three-ones bounds identified
in the research report.  Adding it as an explicitly conditional rule closes
certificate checks for residues $4$ through $8$ over tested ranges of ranks;
those conditional checks are not claimed as exact values.  A uniform proof
of this inverse-height bound, and a general classification for every residue,
remain open.  Hence this submission does not determine $\DB(n)$ for every
$n$ and requests evaluation for partial progress only.

\begin{thebibliography}{9}
\bibitem{GH} J. R. Griggs and C. H. Ho, ``The cycling of partitions and
compositions under repeated shifts,'' \emph{Advances in Applied Mathematics}
21 (1998), 205--227. \url{https://people.math.sc.edu/griggs/cycling.pdf}.
\bibitem{Cell} BAINSA, ``$D_B(n)$ for every $n$,'' P3 C6.
\url{https://hackathon.bainsa.ai/p/p3/c6}.
\end{thebibliography}
\end{document}
